\hypertarget{appendix-b08d3788c5c6}{}
\section*{Return values and copy elision}
\section{Return values and copy elision}\label{return-values-and-copy-elision}

\textbf{Curriculum depth:} Introduced · \textbf{First seen:} C009 · \textbf{Latest development:} C009

\subsection{Definition and mechanism}\label{definition-and-mechanism}

Appropriate same-type prvalues directly initialize result objects since C++17. NRVO for eligible named locals is optional in C++23 and needs a usable fallback. Move-eligible local/parameter return expressions have special xvalue treatment; return std::move(local) loses the ordinary NRVO form.

\subsection{Prerequisites}\label{prerequisites}

\begin{itemize}
\tightlist
\item
  \hyperref[appendix-1121b30eaeb0]{Value categories and materialization} --- The return expression category distinguishes result construction from a bound path or move candidate.
\item
  \hyperref[appendix-a82f9c438857]{Object construction and storage reuse} --- A new destination requires a valid lifetime, independently of available bytes.
\end{itemize}

\subsection{Used by}\label{used-by}

Later sources will extend this node when their guarantees are developed.

\subsection{Recall and next depth}\label{recall-and-next-depth}

Return an immovable type directly, then explain why a named-local return fails. Pending: overlapping-subobject cases and detailed return deduction. Written coverage does not demonstrate personal mastery.

\subsection{Source and navigation}\label{source-and-navigation}

\href{Continuous\%20Notes/Computing\%20Foundations\%20-\%20From\%20Bits\%20to\%20C++\%20and\%20Systems\%20APIs\#Return\%20a\%20result\%20without\%20handing\%20out\%20a\%20lifetime\%20problem}{Teaching mechanism and examples} · \hyperref[appendix-fd090d982a15]{Exact coverage} · \href{}{Dependency Map (02-Foundations-II, appendix-4ad47bdd269a)}
